% ---------------------------------------------------------------------------
% Paper skeleton (from the `paper-writing` skill).
%
% Conventions used here:
%  * Every theorem-like environment carries a comment "% register: R<n>" on the
%    line after \begin{...}, naming its entry in the Results register
%    (PROJECT_SUMMARY.md). The pre-circulation check compares these with the
%    register.
%  * Placeholders are written <<LIKE THIS>>. Search for "<<" before circulating.
%  * Bibliography: BibTeX with amsalpha (or amsplain). We use BibTeX rather
%    than biblatex/biber because arXiv and most journal production systems
%    process BibTeX directly, whereas a biblatex document usually has to ship
%    a .bbl matching the server's package versions.
% ---------------------------------------------------------------------------
\documentclass[11pt]{amsart}

% --- Packages ---------------------------------------------------------------
\usepackage[utf8]{inputenc} % harmless on current engines; needed on old ones
\usepackage[T1]{fontenc}
\usepackage{lmodern}
\usepackage{amssymb}
\usepackage{mathtools}       % loads and extends amsmath
\usepackage{graphicx}
\usepackage[dvipsnames]{xcolor}
% hyperref must come after almost every other package, and cleveref after
% hyperref. Theorem environments are declared after cleveref so it knows each
% environment's name even when environments share a counter.
\usepackage[colorlinks=true, linkcolor=MidnightBlue, citecolor=OliveGreen,
            urlcolor=MidnightBlue]{hyperref}
\usepackage[capitalise,nameinlink]{cleveref}

% --- Theorem environments (numbered within sections, one shared counter) ---
\theoremstyle{plain}
\newtheorem{theorem}{Theorem}[section]
\newtheorem{proposition}[theorem]{Proposition}
\newtheorem{lemma}[theorem]{Lemma}
\newtheorem{corollary}[theorem]{Corollary}
\newtheorem{conjecture}[theorem]{Conjecture}
% Main results in the introduction: Theorem A, Theorem B, ...
\newtheorem{mainthm}{Theorem}
\renewcommand{\themainthm}{\Alph{mainthm}}

\theoremstyle{definition}
\newtheorem{definition}[theorem]{Definition}
\newtheorem{example}[theorem]{Example}
\newtheorem{question}[theorem]{Question}

\theoremstyle{remark}
\newtheorem{remark}[theorem]{Remark}
% Use "observation" (remark style, not plain) for NUMERICAL evidence, so that
% uncertified results never look like theorems.
\newtheorem{observation}[theorem]{Observation}

\crefname{mainthm}{Theorem}{Theorems}
\crefname{conjecture}{Conjecture}{Conjectures}
\crefname{question}{Question}{Questions}
\crefname{observation}{Observation}{Observations}

\numberwithin{equation}{section}

% --- Notation -----------------------------------------------------------------
% Define every symbol once, here (or \input{macros}), and keep it consistent
% with the notation table in PROOF_PROGRESS.md. State conventions that differ
% between sources (signs, normalisations) in Section 2.
\newcommand{\N}{\mathbb{N}}
\newcommand{\Z}{\mathbb{Z}}
\newcommand{\Q}{\mathbb{Q}}
\newcommand{\R}{\mathbb{R}}
\newcommand{\C}{\mathbb{C}}
\DeclareMathOperator{\rank}{rank}
\DeclareMathOperator{\Hom}{Hom}
% <<PROJECT-SPECIFIC MACROS>>

% --- Front matter ---------------------------------------------------------------
\title[<<SHORT TITLE>>]{<<TITLE>>}

\author{<<AUTHOR NAME>>}
\address{<<INSTITUTION AND POSTAL ADDRESS>>}
\email{<<EMAIL>>}

% \subjclass[2020]{Primary <<MSC>>; Secondary <<MSC>>}
% \keywords{<<KEYWORDS>>}

% \date{<<VERSION LABEL, e.g. Working paper v0.1, YYYY-MM-DD>>}

\begin{document}

\begin{abstract}
<<Write last. State the main results at exactly their register status:
conditional results as conditional, numerical evidence as evidence. Nothing here
may be stronger than the body.>>
\end{abstract}

\maketitle

% ===========================================================================
\section{Introduction}\label{sec:intro}

<<Context: the problem, why it matters, what was known (cite exact theorem
numbers).>>

\begin{mainthm}\label{thm:A}
% register: R<<n>>
% Copy the statement from PROOF_PROGRESS.md, edit for readability, then
% compare it hypothesis by hypothesis with the original.
<<Precise statement, all hypotheses included.>>
\end{mainthm}

% A conditional result (register status PROVED-MODULO) is stated with its
% condition as an explicit hypothesis:
% \begin{mainthm}\label{thm:B}
% % register: R<<n>>
% Assume \cref{conj:main}. Then <<...>>.
% \end{mainthm}

<<Method: the idea of the proof in a few paragraphs. What is new compared with
existing approaches.>>

\subsection*{Organisation}
\Cref{sec:prelim} fixes notation and recalls the results we use.
<<...>>

% ===========================================================================
\section{Preliminaries}\label{sec:prelim}

\subsection{Notation and conventions}
<<Notation table in prose. State every convention that differs between
sources, e.g. the sign of the Laplacian or curvature, normalisations.>>

\subsection{Results from the literature}
% State cited results exactly, with theorem numbers, in the version checked.
\begin{theorem}[{<<Author>> \cite[Theorem~<<x.y>>]{<<bibkey>>}}]\label{thm:cited}
% register: R<<n>> (CITED)
<<Exact statement as in the source.>>
\end{theorem}

% ===========================================================================
\section{<<Main argument>>}\label{sec:main}

\begin{lemma}\label{lem:first}
% register: R<<n>>
<<Statement.>>
\end{lemma}

\begin{proof}
<<Keep every justification, including explicit hypothesis checks when
applying cited theorems ("hypothesis (ii) of \cref{thm:cited} holds because
...").>>
\end{proof}

\begin{proof}[Proof of \cref{thm:A}]
<<...>>
\end{proof}

% ===========================================================================
\section{Computations}\label{sec:computations}
% Delete this section if there are none. For each computational claim state:
% what was computed; the method; its certification level (CERTIFIED or
% NUMERICAL); software and versions; where the code, data and run manifest are
% (repository + commit, or archive DOI; run ID).

\begin{proposition}[computer-assisted]\label{prop:certified}
% register: R<<n>> (CERTIFIED)
<<Statement.>>
\end{proposition}

\begin{proof}
<<Describe the certificate (e.g. interval arithmetic, exact rational
arithmetic, certified root isolation) and how to reproduce it.>>
\end{proof}

\begin{observation}\label{obs:numerical}
% register: R<<n>> (NUMERICAL)
<<Numerical evidence, stated as evidence: "In all <<N>> examples computed
(<<scope>>), ...". Never "we verify".>>
\end{observation}

% Figures: generate from scripts in the repository; keep the path next to the
% include.
% \begin{figure}
%   \centering
%   % generated by: <<code/plot_xyz.py>> (run <<YYYY-MM-DD_slug>>)
%   \includegraphics[width=0.7\textwidth]{figures/<<name>>}
%   \caption{<<...>>}\label{fig:<<name>>}
% \end{figure}

% ===========================================================================
\section{Further directions}\label{sec:further}

\begin{conjecture}\label{conj:main}
% register: R<<n>> (CONJECTURE)
<<Statement, followed by a summary of the evidence, labelled by status.>>
\end{conjecture}

\begin{question}\label{q:open}
<<Open questions, including UNRESOLVED register entries and what was tried.>>
\end{question}

% ===========================================================================
\section*{Provenance and status of this manuscript}
<<Specific, accurate account of AI assistance, of who checked what, and of the
manuscript's stage (working paper / preprint / submitted / accepted). See
references/provenance-statements.md in the paper-writing skill. Never describe
a review that did not happen.>>

\section*{Acknowledgements}
<<Funding, people, computing resources.>>

\bibliographystyle{amsalpha}
\bibliography{refs}

\end{document}
